\section{Monitoring Topology}
\label{sec:MonitoringTopology}

\subsection{Purpose}

The monitoring subsystem is a separate runtime concern from the generated process execution graph. The generated handlers emit canonical lifecycle events, while the monitoring topology consumes those events and turns them into queryable operational views.

This separation matters because execution topics are transport-oriented. They show how work is routed. The monitoring topology instead answers questions such as:
\begin{itemize}
\item where is a specific process instance now,
\item how many instances are currently in a given state,
\item which activities are taking the longest,
\item which instances appear stuck.
\end{itemize}

\subsection{Streams topology}

The core implementation lives in \texttt{org.gautelis.durga.monitoring.ProcessMonitoringTopology}. The topology consumes the shared \texttt{process-events} topic, filters out null keys or values, and then groups by process instance id.

From there, the monitoring flow is:
\begin{enumerate}
\item aggregate lifecycle events into a latest per-instance projection,
\item materialize that projection into a state store and a compacted topic,
\item regroup by current state key,
\item count instances per current state,
\item materialize those counts into a second state store and topic.
\end{enumerate}

\begin{figure}[h!]
    \centering
    \begin{tikzpicture}[node distance=15mm,>=latex]
        \tikzstyle{box}=[draw=iptoMuted, rounded corners=2pt, line width=0.8pt, align=center, minimum width=31mm, minimum height=10mm, fill=white]
        \tikzstyle{accentbox}=[box, fill=iptoAccent!12]
        \node[accentbox] (events) {process-events};
        \node[box, right=10mm of events] (aggregate) {aggregate by\\instance id};
        \node[box, right=10mm of aggregate] (state) {process-state\\topic + store};
        \node[box, below=6mm of state] (query) {HTTP + CLI\\query surface};
        \node[box, below=24mm of aggregate] (count) {group by\\current state};
        \node[box, right=10mm of count] (counts) {process-state-counts\\topic + store};

        \draw[->, thick] (events) -- (aggregate);
        \draw[->, thick] (aggregate) -- (state);
        \draw[->, thick] (aggregate) -- (count);
        \draw[->, thick] (count) -- (counts);
        \draw[->, thick] (state) -- (query);
        \draw[->, thick] (counts) -- (query);
    \end{tikzpicture}
    \caption{Durga monitoring topology from canonical lifecycle stream to queryable state}
    \label{fig:DurgaMonitoringTopology}
\end{figure}

\subsection{Materialized state}

The per-instance projection is represented by \texttt{ProcessStateView}. It is not just a copy of the latest raw event. It accumulates enough information to support operational queries, including:
\begin{itemize}
\item current activity or lifecycle state,
\item started, updated and completed timestamps,
\item retry and failure information,
\item per-activity timing information,
\item correlation and business-key metadata.
\end{itemize}

That projection is written both to the Kafka Streams state store and to the \texttt{process-state} topic. The state store supports local low-latency queries. The topic keeps the projection observable and externalizable.

\subsection{Counts, latency and stuck detection}

The counts projection is simpler. The topology derives a state key from the latest per-instance view and counts how many instances currently map to each such key.

Latency and stuck-instance reporting are handled by the query service layer on top of the materialized state:
\begin{itemize}
\item latency summaries are computed from timestamps and durations recorded in \texttt{ProcessStateView},
\item stuck-instance detection scans state entries and compares their age against a threshold,
\item process-specific counts are filtered from the global count view.
\end{itemize}

\subsection{HTTP and CLI surface}

The query service is exposed through \texttt{ProcessMonitoringHttpServer} and the companion CLI client. The main endpoints are:
\begin{itemize}
\item \texttt{/health},
\item \texttt{/instances/\{processInstanceId\}},
\item \texttt{/processes/\{processId\}/counts},
\item \texttt{/processes/\{processId\}/latency},
\item \texttt{/counts},
\item \texttt{/stuck}.
\end{itemize}

This is complemented by a lightweight dashboard page that polls the same query surface.

\subsection{Relationship to demos}

The monitoring topology is not only useful when a generated BPMN runtime is running. Repository-level demo publishers and scenario runners can emit synthetic lifecycle streams directly to \texttt{process-events}. That makes the monitoring subsystem independently testable and also makes it easier to understand the read-model behavior before a full generated process is involved.

\subsection{Current limits}

The topology is intentionally direct. It is strong enough for local experimentation and architectural validation, but it is not yet a comprehensive production operations subsystem. The main current limits are:
\begin{itemize}
\item the read model is local to the monitoring app instance,
\item some queries are derived by scanning materialized state rather than by maintaining separate specialized indexes,
\item history and trend reporting remain lighter than current-state reporting.
\end{itemize}

Even with those limits, the monitoring path is a major part of what makes Durga usable as an exploration tool rather than only a code generator.
